\uuid{6VKd}
\exo7id{2325}
\titre{exo7 2325}
\auteur{matexo1}
\organisation{exo7}
\datecreate{2002-02-01}
\isIndication{false}
\isCorrection{false}
\chapitre{Calcul d'intégrales}
\sousChapitre{Primitives diverses}
\cpge{mpsi-mp2i.primitives-equations-differentielles.a, pcsi.techniques-calcul-primitives.a, ptsi.techniques-calcul-primitives.a, mpsi-mp2i.fonctions-valeurs-reelles-complexes.c, pcsi.techniques-calcul-fonctions.c, ptsi.techniques-calcul-fonctions.c, bcpst1.derivees-primitives-integrales.c}
\module{Analyse}
\niveau{L1}
\difficulte{}

\contenu{
\question{
Calculer les primitives des fonctions suivantes\,:
\begin{eqnarray*}
t &\mapsto & \tan t, \\
t &\mapsto & \arg \sinh t, \\
t &\mapsto & {\tan^3 t \over {\cos^6 t }}, \\
t &\mapsto & {1 \over {\sqrt{t} + \sqrt[3]{t}}}, \\
t &\mapsto & {1 \over {2\sqrt {t} + \sqrt{t+2}}}, \\
t &\mapsto & t \sqrt[n]{1+t}, \\
t &\mapsto & \cosh^3 t, \\
t &\mapsto & {t^3 \over {(a^2-t^2)^{3/2} }}, \\
t &\mapsto & {\sqrt[3]{1+\sqrt[4]t} \over {\sqrt{t}}}, \\
t &\mapsto & \sqrt{t \sqrt{t \sqrt{t}}}, \\
t &\mapsto & { {1 \over {t^2}} \sqrt[4]{t+1 \over {t-1}}}.
\end{eqnarray*}
}
}
